\hypertarget{chapter-06-page-127}{}
\subsection{原子 $H^1$ 空间}
\label{sec:ch6-atomic-h1}

在第~\ref{sec:ch3-riesz-kolmogorov-theorems}~节中, 我们看到可积函数的 Hilbert 变换一般不属于 $L^1$, 而函数积分为零是其属于 $L^1$ 的必要条件. 这里要定义 $L^1$ 的一个子空间, 使任意奇异积分作用于其中的函数后都属于 $L^1$. 我们先定义它的基本元素---原子.

\begin{Definition}
\label{def:ch6-atom}
原子是定义在 $\mathbb R^n$ 上的复值函数 $a$, 它支于某个立方体 $Q$, 并且
\[
 \int_Qa(x)\,dx=0,
 \qquad
 \|a\|_\infty\le\frac1{|Q|}.
\]
\end{Definition}

\begin{Proposition}
\label{prop:ch6-singular-integral-on-atoms}
设 $T$ 满足定理~\ref{thm:ch5-generalized-calderon-zygmund-bounds} 的假设, 且其核满足条件 \eqref{eq:ch5-generalized-hormander-y}. 则存在常数 $C$, 使每个原子 $a$ 都满足
\[
 \|Ta\|_1\le C.
\]
\end{Proposition}

\begin{Proof}
因为 $a\in L^2$, 所以 $Ta$ 有定义. 设 $Q^*$ 与 $Q$ 有相同的中心 $c_Q$, 边长是 $Q$ 的 $2\sqrt n$ 倍. 由 $T$ 在 $L^2$ 上有界,
\[
 \int_{Q^*}|Ta(x)|\,dx
 \le |Q^*|^{1/2}\left(\int_{Q^*}|Ta(x)|^2\,dx\right)^{1/2}
 \le C|Q|^{1/2}\left(\int_Q|a(x)|^2\,dx\right)^{1/2}
 \le C.
\]
此外, 由 $a$ 的平均值为零以及条件 \eqref{eq:ch5-generalized-hormander-y},
\[
\begin{aligned}
 \int_{\mathbb R^n\setminus Q^*}|Ta(x)|\,dx
 &=\int_{\mathbb R^n\setminus Q^*}
   \left|\int_QK(x,y)a(y)\,dy\right|dx\\
 &=\int_{\mathbb R^n\setminus Q^*}
   \left|\int_Q[K(x,y)-K(x,c_Q)]a(y)\,dy\right|dx\\
 &\le\int_Q\int_{\mathbb R^n\setminus Q^*}
   |K(x,y)-K(x,c_Q)|\,dx\,|a(y)|\,dy
 \le C.
\end{aligned}
\]
\end{Proof}

\hypertarget{chapter-06-page-128}{}
现在定义原子 $H^1$ 空间, 记作 $H^1_{\mathrm{at}}$, 令
\[
 H^1_{\mathrm{at}}(\mathbb R^n)
 =\left\{\sum_j\lambda_ja_j:
 a_j\text{ 是原子},\ \lambda_j\in\mathbb C,\ \sum_j|\lambda_j|<\infty\right\}.
\]
显然 $H^1_{\mathrm{at}}\subset L^1$. 在 $H^1_{\mathrm{at}}$ 上定义范数
\[
 \|f\|_{H^1_{\mathrm{at}}}
 =\inf\left\{\sum_j|\lambda_j|:f=\sum_j\lambda_ja_j\right\}.
\]
关于此范数, $H^1_{\mathrm{at}}$ 是 Banach 空间. 细节留给读者. 由命题~\ref{prop:ch6-singular-integral-on-atoms} 立即可知, 奇异积分从 $H^1_{\mathrm{at}}$ 到 $L^1$ 有界.

\begin{Corollary}
\label{cor:ch6-singular-integral-h1-to-l1}
设 $T$ 是命题~\ref{prop:ch6-singular-integral-on-atoms} 中的算子, 且 $f\in H^1_{\mathrm{at}}$. 则
\[
 \|Tf\|_1\le C\|f\|_{H^1_{\mathrm{at}}}.
\]
\end{Corollary}

在如下意义下, $H^1_{\mathrm{at}}$ 是使推论~\ref{cor:ch6-singular-integral-h1-to-l1} 成立的 $L^1(\mathbb R^n)$ 的最大子空间. 设 $R_1,\ldots,R_n$ 是 $\mathbb R^n$ 上的 Riesz 变换, 当 $n=1$ 时就是 Hilbert 变换, 并定义
\[
 H^1(\mathbb R^n)
 =\{f\in L^1(\mathbb R^n):R_jf\in L^1(\mathbb R^n),\ 1\le j\le n\},
\]
其范数为
\[
 \|f\|_{H^1}=\|f\|_1+\sum_{j=1}^n\|R_jf\|_1.
\]

\begin{Theorem}
\label{thm:ch6-h1-atomic-equivalence}
$H^1(\mathbb R^n)=H^1_{\mathrm{at}}(\mathbb R^n)$, 且二者的范数等价.
\end{Theorem}

这里不证明定理~\ref{thm:ch6-h1-atomic-equivalence}; 证明可见 García-Cuerva 与 Rubio de Francia \MBUnresolvedReference{citation}{[6, Chapter 3]}, 或 Stein \MBUnresolvedReference{citation}{[17, Chapter 3]}.
